% ПРИЛОЖЕНИЕ, без номер. Само указател към изходния текст, без вмъкване на код.
\section*{ПРИЛОЖЕНИЕ}
\label{sec:appendix}
\addcontentsline{toc}{section}{ПРИЛОЖЕНИЕ}

Изходният текст не се препечатва тук. Хранилището е публично:
\url{https://github.com/Alex-Tsvetanov/conduit}. Табл.~\ref{tab:sources} дава указател към
него; редовете са преброени с \code{wc -l}.

\begin{table}[H]
\centering
\caption{Указател към изходния текст на проекта}
\label{tab:sources}
\begin{tabular}{L{4.4cm}L{7.6cm}L{1.6cm}}
\toprule
\textbf{Път} & \textbf{Какво съдържа} & \textbf{Редове} \\
\midrule
\code{bytes.hpp}, \code{task.hpp}, \code{net.*}, \code{tls.*} & буфери, \code{task<T>}, сокет, TLS и цикъл на събитията & 1509 \\
\code{crypto.*}, \code{scram.*} & MD5, SHA-1, SHA-256, HMAC, PBKDF2, SCRAM-SHA-256 & 470 \\
\code{codec.hpp}, \code{value.*}, \code{trace.*} & договорът между кодеците, декодерите, броене на съобщенията & 612 \\
\code{lru\_cache.hpp}, \code{row\_map.hpp}, \code{pool.hpp} & кеш с изместване, типово изображение, асинхронен пул & 496 \\
\code{pg/}, \code{src/pg\_*.cpp} & кодек и връзка за PostgreSQL, двата цикъла, TLS & 1176 \\
\code{mysql/}, \code{src/mysql\_*.cpp} & кодек и връзка за MySQL, текст протокол, TLS & 852 \\
\midrule
\code{tests/} & 91 модулни и 14 интеграционни случая, собствена рамка & 2406 \\
\code{benchmarks/bench.cpp} & измервателната програма, флаг \code{--raw} & 613 \\
\code{benchmarks/symmetric\_transfer\_check.cpp} & проверка дали симетричният трансфер расте по стек & 85 \\
\code{examples/demo.cpp} & показна програма с отпечатан протоколен обмен & 286 \\
\code{docker-compose.yml} & PostgreSQL 16.4 и MySQL 8.0.39, фиксирани версии, SSL & 87 \\
\bottomrule
\end{tabular}
\end{table}

Схемата на измервателната таблица се създава от самата измервателна програма и е част от
нейния изходен текст. Суровите проби от всяко повторение се записват при подаден флаг
\code{--raw} във файла \code{raw\_samples.csv}; пускът, от който са таблиците в
Раздел~\ref{sec:results}, произведе 5660 записа с полета \code{case}, \code{repetition} и
\code{microseconds}.
